\section{International Transparency-Indexed Fiat Framework (ITIF)}
\subsection{Purpose and Scope}
ITIF is a transparency-first interpretive regime for sovereign fiat systems. It is not a new currency or supranational authority. It preserves national sovereignty while introducing a common language for fiscal space and monetary integrity.

\subsection{Core Design}
\begin{itemize}
  \item Transparency overlay: standardized disclosure of fiscal space, bond structure resilience, and tax clarity.
  \item T-value: a composite integrity band (T1--T10) reflecting audit transparency, fiscal space elasticity, data integrity (CLEAR verification), bond resilience, tax leakage, and CB--treasury synchronization.
  \item Integration: interoperable with IMF Article IV, BIS/FSB data schemas, and G20 transparency codes.
\end{itemize}

\subsection{Governance and Safeguards}
\begin{itemize}
  \item Operated by national authorities; no delegation of issuance power.
  \item Public T-value bands to anchor market pricing and reduce opacity premia.
  \item CLEAR provides a non-intrusive verification layer; it does not store personal data, does not perform inference, and fully preserves national sovereignty.
\end{itemize}

\subsection{Roadmap (Condensed)}
\begin{enumerate}
  \item Transparency baseline: publish T-value components; enable CLEAR-based audit trails for pilots.
  \item Computation layer: national statistics offices compute provisional T-values; multilateral observers read-only.
  \item Sovereign compatibility: tag bonds with T-value metadata; markets price transparency-adjusted risk.
  \item International integration: coordinated publication cycle; eligibility criteria for transparency-linked channels.
  \item Optional ledger convergence: ITIF ledger for cross-border reconciliation with domestic control retained.
\end{enumerate}

\subsection{Reference Diagrams}
\placeholderfigure{Tvalue\_Structure.pdf}{T-value components}
\placeholderfigure{Policy\_Stack\_Diagram.pdf}{Ultimate--ITIF--UBI-T--CLEAR stack}
